\documentclass{zhvt-classic}

% 使用本地源流明体（GenRyuMin2 TW）按文件路径加载，不依赖系统字体名
% 版框（内框／界栏／书口侧栏／鱼尾）改用与版心印章同一印泥棕色：
%   印章图片 yshs.png 实色 = RGB(135,68,52) = HTML 874434
\usepackage{xcolor}
\usepackage{tikz}
% sealbrown 双注册：xcolor 供 TikZ/pgf，L3 供 cls 版框
\definecolor{sealbrown}{HTML}{874434}
\ExplSyntaxOn
\color_set:nnn { sealbrown } { HTML } { 874434 }
\ExplSyntaxOff

\zhvtset{
  grid color       = sealbrown,
  main font        = GenRyuMin2TW-SB,
  main font path   = C:/Users/Administrator/texmf/fonts/opentype/public/GenRyuMin2TW/,
jiazhu font      = GenRyuMin2TW-M,
  jiazhu font path = C:/Users/Administrator/texmf/fonts/opentype/public/GenRyuMin2TW/,
  % 页码框容纳 4 字：正常页「九十七」+前置空 1 格=4 行恰好；
  % 100+ 页「一百零一」共 5 行超出，末字「一/二/三」被框裁掉（2026-10-04 实测），故放宽到 5
  shukou page number chars = 5,
  % 书口书名框高从默认 6 字减到 4 字，把省出的 2 字高拨给 chapter 框，
  % 使 6 字书名+卷名放进 9 字高的 chapter 框后不再压线（2026-10-05）
  shukou title chars = 4,
}

\title{史记}[孝武本纪]
\author{汉·司马迁}
\maker{闲叟漫记·重製}


% ===== 索隐述赞样式（2026-11-08 四次定稿 · 愚大人亲手改定参数）=====
% ⚠ 竖排前提：自定义字体家族必须带 RawFeature={vertical:+vert}；
%   simhei.ttf 加 +vhal 会整体加载失败静默回退，只写 +vert；simfang 同理。
\newCJKfontfamily{\shuzanfont}[Path=C:/Windows/Fonts/,RawFeature={vertical:+vert}]{simhei.ttf}
\newCJKfontfamily{\shuzanbody}[Path=C:/Windows/Fonts/,RawFeature={vertical:+vert}]{simfang.ttf}
% 题头：152x30pt 圆角 6pt 棕底白字；baseline (0,2) 使整框落入界行格子不压版框线；
% 白字节点 x 不补偿、y +8.2pt（像素级墨迹实测 dx/dy≈0，L/R 留白约 20pt）。
\newcommand{\shuzanhead}[1]{%
  \begin{tikzpicture}[baseline={(0,2)}]
    \fill[sealbrown, rounded corners=6pt]
      (-76pt,-15pt) rectangle (76pt,15pt);
    \node[text=white, inner sep=0pt, anchor=center] at (0pt,8.2pt)
      {\shuzanfont\fontsize{23}{23}\selectfont #1};
  \end{tikzpicture}%
}
% ★ 修正书口卷名位置：chapter 框起始位置改回 cls 默认值，避免全局影响。
%   只在每个卷首页通过 \juan 命令里的空命令微调该页位置（安全，不会跨页污染）。
\ExplSyntaxOn
\AtBeginDocument{
  \dim_set:Nn \l__zhvt_shukou_chapter_title_start_dim
    { \l__zhvt_shukou_title_height_dim + \l__zhvt_font_size_dim * 2 }
  \dim_set:Nn \l__zhvt_shukou_page_number_start_dim
    { \l__zhvt_shukou_chapter_title_start_dim + \l__zhvt_shukou_chapter_title_height_dim }
}
\ExplSyntaxOff

\begin{document}

% ★★ 版心（书口）排布覆写：把书名放到鱼尾之下（2026-10-04）
%
%   cls 默认版心自上而下：书名（title 静态框，在鱼尾之上）→ 鱼尾 → 篇名（chapter 动态框）→ 页码
%   本模板改为：          鱼尾（位置不动）→ 书名 → 篇名 → 页码分割线 → 页码
%
%   手法：把 cls 放在 title 静态框里的书名，改放进 chapter 动态框（鱼尾之下）；
%         title 静态框留空。书名取 \l__zhvt_title_str，即 \title{...} 的第一参数，
%         所以换书时只改 \title，不必在此重复写书名。
%
%   ⚠ 三条必须遵守的写法（2026-10-04 实测）：
%     ① 汉字必须用 \hbox{...} 包住。flowfram 会把框内容装进一个宽 14.4mm（= 一个 line_sep
%        = 书口栏宽）的纵向 minipage 并按该宽度折行，而整页靠 PDF /Rotate 90 显示、
%        框内文字本身不变向 —— 不包 \hbox 时「围炉夜话」4 字会被逐字折成 4 行，
%        在版面上就成了横排 4 字并横向溢出版框。
%     ② 竖排下**一个 \hbox 就是一整列**。书名与篇名要同处一列，必须放进同一个 \hbox；
%        分成两个 \hbox 时，后一个会被当成新的一列、甩到隔壁栏去
%        （实测「序」跑到书口栏外侧 14.4mm 的正文列上）。中间的 \quad 在框内是横向空白，
%        显示出来正好是书名与篇名之间的纵向间隔。
%     ③ ★ 内容盒要写成 **与框等宽的定宽盒** `\hbox to \l__zhvt_shukou_chapter_title_height_dim{…\hfil}`。
%        实测：chapter 框是把内容盒沿纵向**居中**摆放的，内容宽度一变，书名就跟着漂：
%        自然宽度的 \hbox（书名 38.4mm）落在 x=157.2mm，加篇名后（57.6mm）落在 x=152.4mm，
%        目录页（无篇名）与正文页的书名就差半格。做成等宽盒后，无论有没有篇名，
%        内容固定从框顶起排，实测两页均落在 x=142.8mm，不再漂。
\ExplSyntaxOn
% 把书口 chapter 框（书名+篇名）起始位置下移 1 字高，避免 4 字卷名顶到鱼尾
% ⚠ cls 的 \__zhvt_calculate_dimensions: 在 \AtBeginDocument 里执行，会重算该 dim；
%   必须把覆写挂到 begin document 之后，否则会被二次覆盖（2026-10-05 实测）。
% (旧方案已移到导言区，直接修改 chapter 框起始位置)
% 卷标题命令 #1=书口篇名(鱼尾) #2=正文卷名（含卷序，卷名内空格用 ~）
%
% ★★ 目录页两项修订（2026-10-05，愚大人指出）★★
%   ① 页码起点统一对齐：
%      原状 —— 目录项写成 \par\hspace{4em}卷名\hspace{4em}页码，页码起点 = 卷名实宽 + 4em。
%      卷名实宽随字数变：「卷一~梁惠王上」6 字 +~、「卷七~离娄上」5 字 +~，
%      差 1 个字位 → 卷七~卷十 4 条的页码整体比其余 10 条高 27.21pt（一个字），看着上下不齐。
%      改法 —— 卷名补全角空格（宽 = 1 字 = 27.215pt）补足到 **7 字符位**再写进目录，
%      14 条页码起点遂全部落在同一条竖线上（起点对齐；末字随页码字数自然参差）。
%      ⚠ 只对目录里的副本 \l_tmpc_tl 补，正文卷名仍用 \l_tmpa_tl 原样 —— 正文像素不受影响。
%   ② 卷名整体上移半字位：前导 \hspace 由 4em 改 3.5em。
%      实测原状目录块在版框内偏下（天头留白 139.6pt / 地脚留白 110.7pt），
%      上移 0.5em = 13.6pt 后上下基本均分，视觉居中。
\NewDocumentCommand { \juan } { mm } {
	\clearpage
	\tl_set:Nn \l__zhvt_chapter_title_tl { #1 }
	\tl_set:Nn \l_tmpa_tl { #2 }
	\tl_set_eq:NN \l_tmpc_tl \l_tmpa_tl
	\int_step_inline:nn { 7 - \str_count:n { #2 } }
	  { \tl_put_right:Nn \l_tmpc_tl { 　 } }
	\tl_set:Nx \l_tmpb_tl { \zhnum { page } }
	\tl_put_left:Nn \l_tmpb_tl { \hspace { 4em } }
	\tl_put_left:Nn \l_tmpb_tl { \l_tmpc_tl }
	\tl_put_left:Nn \l_tmpb_tl { \par \hspace { 3.5em } }
	\addtocontents { toc } { \l_tmpb_tl }
	\skip_horizontal:n { \l__zhvt_font_size_dim * 4 }
	\l_tmpa_tl
	\par
	\thispagestyle { empty }
}
\cs_set_protected:Npn \__zhvt_set_shukou_content: {
	% ★ 鱼尾上方留白（title 静态框）嵌入闲叟钤印：
	%   整页 /Rotate 90，故此处把图用 \rotatebox{90} 转正为竖章；
	%   尺寸卡在版心条幅内（书内约 7mm 宽 × 33.8mm 高，留余量不压鱼尾）。
	\setstaticcontents * { title } {
		\vbox_to_ht:nn { \l__zhvt_baseline_skip_dim } {
			\vskip 14.2pt
			\hbox_to_wd:nn { \l__zhvt_shukou_title_height_dim } {
				\hfil
				\rotatebox [ origin = c ] { 90 } { \includegraphics [ width = 7mm , height = 33.8mm ] { yshs.png } }
				\hfil
			}
		}
	}
% ★ 版心（书口）统一改法（2026-11-16 愚大人定）：18pt＋框内居中＋书口文字同版框色
%   ① 原 27.2pt 时 7 字长篇名（如「廉颇蔺相如列传」）末字溢出 chapter 框、压住页码；
%      改 \fontsize{18pt}{18pt}，前后 \hfil 令字串在框内居中（不随篇名字数漂移）。
%   ② \vbox{\vskip 8.2pt} 定量下移，抵消缩字号后栏宽方向 8.2pt 上偏（像素实测 y 中心 543.3 = 原版）。
%      ⚠ \raisebox / \vspace* 在 flowfram 框内无效；换字号须同步换补偿：16pt→10.1、19pt→7.4。
%   ③ chapter 与 page~number 均置 \color{sealbrown}（HTML 874434），与版框／界行同色；正文仍黑。
%   ⚠ 页码字位、钤印、鱼尾、界行、正文一律未动；五部合订本同此口径。
	\setdynamiccontents * { chapter }
	{ \color{sealbrown} \vbox { \vskip 8.2pt \hbox to \l__zhvt_shukou_chapter_title_height_dim
		{ \hfil \fontsize{18pt}{18pt}\selectfont \l__zhvt_title_str \quad \l__zhvt_chapter_title_tl \hfil } } }
	\setdynamiccontents * { page~number } { \color{sealbrown} \quad \zhnum { page } }
}
\__zhvt_set_shukou_content:
\ExplSyntaxOff

\maketitle

% 单篇，无多卷目录，不渲染目录页

\juan{孝武本纪}{卷十二~孝武本纪第十二}
\hspace*{2em}孝武皇帝者。孝景中子也。母曰王太后。孝景四年。以皇子为胶东王。孝景七年。栗太子废为临江王。以胶东王为太子。孝景十六年崩。太子即位。为孝武皇帝。孝武皇帝初即位。尤敬鬼神之祀。\\
\hspace*{2em}元年。汉兴已六十余岁矣。天下泰安。荐绅之属皆望天子封禅改正度也。而上乡儒术。招贤良。赵绾。王臧等以文学为公卿。欲议古立明堂城南。以朝诸侯。草巡狩封禅改历服色事未就。会窦太后治黄老言。不好儒术。使人微得赵绾等奸利事。召案绾。臧。绾。臧自杀。诸所兴为者皆废。\\
\hspace*{2em}后六年。窦太后崩。其明年。上征文学之士公孙弘等。\\
\hspace*{2em}明年。上初至雍。郊见五畤。后常三岁一郊。是时上求神君。舍之上林中蹄氏观。神君者。长陵女子。以子死悲哀。故见神于先后宛若。宛若祠之其室。民多往祠。平原君往祠。其后子孙以尊显。及武帝即位。则厚礼置祠之内中。闻其言。不见其人云。\\
\hspace*{2em}是时而李少君亦以祠灶。谷道。却老方见上。上尊之。少君者。故深泽侯入以主方。匿其年及所生长。常自谓七十。能使物。却老。其游以方遍诸侯。无妻子。人闻其能使物及不死。更馈遗之。常余金钱帛衣食。人皆以为不治产业而饶给。又不知其何所人。愈信。争事之。少君资好方。善为巧发奇中。尝从武安侯饮。坐中有年九十余老人。少君乃言与其大父游射处。老人为儿时从其大父行。识其处。一坐尽惊。少君见上。上有故铜器。问少君。少君曰。﹁此器齐桓公十年陈于柏寝。﹂已而案其刻。果齐桓公器。一宫尽骇。以少君为神。数百岁人也。\\
\hspace*{2em}少君言于上曰。﹁祠灶则致物。致物而丹沙可化为黄金。黄金成以为饮食器则益寿。益寿而海中蓬莱仙者可见。见之以封禅则不死。黄帝是也。臣尝游海上。见安期生。食臣枣。大如瓜。安期生仙者。通蓬莱中。合则见人。不合则隐。﹂于是天子始亲祠灶。而遣方士入海求蓬莱安期生之属。而事化丹沙诸药齐为黄金矣。\\
\hspace*{2em}居久之。李少君病死。天子以为化去不死也。而使黄锤史宽舒受其方。求蓬莱安期生莫能得。而海上燕齐怪迂之方士多相效。更言神事矣。\\
\hspace*{2em}亳人薄诱忌奏祠泰一方。曰。﹁天神贵者泰一。泰一佐曰五帝。古者天子以春秋祭泰一东南郊。用太牢具。七日。为坛开八通之鬼道。﹂于是天子令太祝立其祠长安东南郊。常奉祠如忌方。其后人有上书。言﹁古者天子三年一用太牢具祠神三一。天一。地一。泰一﹂。天子许之。令太祝领祠之忌泰一坛上。如其方。后人复有上书。言﹁古者天子常以春秋解祠。祠黄帝用一枭破镜。冥羊用羊。祠马行用一青牡马。泰一。皋山山君。地长用牛。武夷君用干鱼。阴阳使者以一牛﹂。令祠官领之如其方。而祠于忌泰一坛旁。\\
\hspace*{2em}其后。天子苑有白鹿。以其皮为币。以发瑞应。造白金焉。\\
\hspace*{2em}其明年。郊雍。获一角兽。若麃然。有司曰。﹁陛下肃祗郊祀。上帝报享。锡一角兽。盖麟云。﹂于是以荐五畤。畤加一牛以燎。赐诸侯白金。以风符应合于天地。\\
\hspace*{2em}于是济北王以为天子且封禅。乃上书献泰山及其旁邑。天子受之。更以他县偿之。常山王有罪。迁。天子封其弟于真定。以续先王祀。而以常山为郡。然后五岳皆在天子之郡。\\
\hspace*{2em}其明年。齐人少翁以鬼神方见上。上有所幸王夫人。夫人卒。少翁以方术盖夜致王夫人及灶鬼之貌云。天子自帷中望见焉。于是乃拜少翁为文成将军。赏赐甚多。以客礼礼之。文成言曰。﹁上即欲与神通。宫室被服不象神。神物不至。﹂乃作画云气车。及各以胜日驾车辟恶鬼。又作甘泉宫。中为台室。画天。地。泰一诸神。而置祭具以致天神。居岁余。其方益衰。神不至。乃为帛书以饭牛。详弗知也。言此牛腹中有奇。杀而视之。得书。书言甚怪。天子疑之。有识其手书。问之人。果伪书。于是诛文成将军而隐之。\\
\hspace*{2em}其后则又作柏梁。铜柱。承露仙人掌之属矣。\\
\hspace*{2em}文成死明年。天子病鼎湖甚。巫医无所不致。不愈。游水发根乃言曰。﹁上郡有巫。病而鬼下之。﹂上召置祠之甘泉。及病。使人问神君。神君言曰。﹁天子毋忧病。病少愈。强与我会甘泉。﹂于是病愈。遂幸甘泉。病良已。大赦天下。置寿宫神君。神君最贵者太一。其佐曰大禁。司命之属。皆从之。非可得见。闻其音。与人言等。时去时来。来则风肃然也。居室帷中。时昼言。然常以夜。天子祓。然后入。因巫为主人。关饮食。所欲者言行下。又置寿宫。北宫。张羽旗。设供具。以礼神君。神君所言。上使人受书其言。命之曰﹁画法﹂。其所语。世俗之所知也。毋绝殊者。而天子独喜。其事秘。世莫知也。\\
\hspace*{2em}其后三年。有司言元宜以天瑞命。不宜以一二数。一元曰建元。二元以长星曰元光。三元以郊得一角兽曰元狩云。\\
\hspace*{2em}其明年冬。天子郊雍。议曰。﹁今上帝朕亲郊。而后土毋祀。则礼不答也。﹂有司与太史公。祠官宽舒等议。﹁天地牲角茧栗。今陛下亲祀后土。后土宜于泽中圜丘为五坛。坛一黄犊太牢具。已祠尽瘗。而从祠衣上黄。﹂于是天子遂东。始立后土祠汾阴脽上。如宽舒等议。上亲望拜。如上帝礼。礼毕。天子遂至荥阳而还。过雒阳。下诏曰。﹁三代邈绝。远矣难存。其以三十里地封周后为周子南君。以奉先王祀焉。﹂是岁。天子始巡郡县。侵寻于泰山矣。\\
\hspace*{2em}其春。乐成侯上书言栾大。栾大。胶东宫人。故尝与文成将军同师。已而为胶东王尚方。而乐成侯姊为康王后。毋子。康王死。他姬子立为王。而康后有淫行。与王不相中。相危以法。康后闻文成已死。而欲自媚于上。乃遣栾大因乐成侯求见言方。天子既诛文成。后悔恨其早死。惜其方不尽。及见栾大。大悦。大为人长美。言多方略。而敢为大言。处之不疑。大言曰。﹁臣尝往来海中。见安期。羡门之属。顾以为臣贱。不信臣。又以为康王诸侯耳。不足予方。臣数言康王。康王又不用臣。臣之师曰。﹃黄金可成。而河决可塞。不死之药可得。仙人可致也。﹄臣恐效文成。则方士皆掩口。恶敢言方哉。﹂上曰。﹁文成食马肝死耳。子诚能修其方。我何爱乎。﹂大曰。﹁臣师非有求人。人者求之。陛下必欲致之。则贵其使者。令有亲属。以客礼待之。勿卑。使各佩其信印。乃可使通言于神人。神人尚肯邪不邪。致尊其使。然后可致也。﹂于是上使先验小方。斗旗。旗自相触击。\\
\hspace*{2em}是时上方忧河决。而黄金不就。乃拜大为五利将军。居月余。得四金印。佩天士将军。地士将军。大通将军。天道将军印。制诏御史。﹁昔禹疏九江。决四渎。间者河溢皋陆。堤繇不息。朕临天下二十有八年。天若遗朕士而大通焉。干称﹃蜚龙﹄。﹃鸿渐于般﹄。意庶几与焉。其以二千户封地士将军大为乐通侯。﹂赐列侯甲第。僮千人。乘舆斥车马帷帐器物以充其家。又以卫长公主妻之。赍金万斤。更名其邑曰当利公主。天子亲如五利之第。使者存问所给。连属于道。自大主将相以下。皆置酒其家。献遗之。于是天子又刻玉印曰﹁天道将军﹂。使使衣羽衣。夜立白茅上。五利将军亦衣羽衣。立白茅上受印。以示弗臣也。而佩﹁天道﹂者。且为天子道天神也。于是五利常夜祠其家。欲以下神。神未至而百鬼集矣。然颇能使之。其后治装行。东入海。求其师云。大见数月。佩六印。贵振天下。而海上燕齐之间。莫不扼捥而自言有禁方。能神仙矣。\\
\hspace*{2em}其夏六月中。汾阴巫锦为民祠魏脽后土营旁。见地如钩状。掊视得鼎。鼎大异于众鼎。文镂毋款识。怪之。言吏。吏告河东太守胜。胜以闻。天子使使验问巫锦得鼎无奸诈。乃以礼祠。迎鼎至甘泉。从行。上荐之。至中山。晏温。有黄云盖焉。有麃过。上自射之。因以祭云。至长安。公卿大夫皆议请尊宝鼎。天子曰。﹁间者河溢。岁数不登。故巡祭后土。祈为百姓育谷。今年丰庑未有报。鼎曷为出哉。﹂有司皆曰。﹁闻昔大帝兴神鼎一。一者一统。天地万物所系终也。黄帝作宝鼎三。象天地人也。禹收九牧之金。铸九鼎。皆尝鬺烹上帝鬼神。遭圣则兴。迁于夏商。周德衰。宋之社亡。鼎乃沦伏而不见。颂云。﹃自堂徂基。自羊徂牛。鼐鼎及鼒。不虞不骜。胡考之休。﹄今鼎至甘泉。光润龙变。承休无疆。合兹中山。有黄白云降盖。若兽为符。路弓乘矢。集获坛下。报祠大飨。惟受命而帝者心知其意而合德焉。鼎宜见于祖祢。藏于帝廷。以合明应。﹂制曰。﹁可。﹂\\
\hspace*{2em}入海求蓬莱者。言蓬莱不远。而不能至者。殆不见其气。上乃遣望气佐候其气云。\\
\hspace*{2em}其秋。上幸雍。且郊。或曰。﹁五帝。泰一之佐也。宜立泰一而上亲郊之。﹂上疑未定。齐人公孙卿曰。﹁今年得宝鼎。其冬辛巳朔旦冬至。与黄帝时等。﹂卿有札书曰。﹁黄帝得宝鼎宛朐。问于鬼臾区。区对曰。﹃帝得宝鼎神䇲。是岁己酉朔旦冬至。得天之纪。终而复始。﹄于是黄帝迎日推䇲。后率二十岁得朔旦冬至。凡二十推。三百八十年。黄帝仙登于天。﹂卿因所忠欲奏之。所忠视其书不经。疑其妄书。谢曰。﹁宝鼎事已决矣。尚何以为。﹂卿因嬖人奏之。上大说。召问卿。对曰。﹁受此书申功。申功已死。﹂上曰。﹁申功何人也。﹂卿曰。﹁申功。齐人也。与安期生通。受黄帝言。无书。独有此鼎书。曰﹃汉兴复当黄帝之时。汉之圣者在高祖之孙且曾孙也。宝鼎出而与神通。封禅。封禅七十二王。唯黄帝得上泰山封﹄。申功曰。﹃汉主亦当上封。上封则能仙登天矣。黄帝时万诸侯。而神灵之封居七千。天下名山八。而三在蛮夷。五在中国。中国华山。首山。太室。泰山。东莱。此五山黄帝之所常游。与神会。黄帝且战且学仙。患百姓非其道。乃断斩非鬼神者。百余岁然后得与神通。黄帝郊雍上帝。宿三月。鬼臾区号大鸿。死葬雍。故鸿冢是也。其后黄帝接万灵明廷。明廷者。甘泉也。所谓寒门者。谷口也。黄帝采首山铜。铸鼎于荆山下。鼎既成。有龙垂胡髯下迎黄帝。黄帝上骑。群臣后宫从上龙七十余人。龙乃上去。余小臣不得上。乃悉持龙髯。龙髯拔。堕黄帝之弓。百姓仰望黄帝既上天。乃抱其弓与龙胡髯号。故后世因名其处曰鼎湖。其弓曰乌号。﹄﹂于是天子曰。﹁嗟乎。吾诚得如黄帝。吾视去妻子如脱躧耳。﹂乃拜卿为郎。东使候神于太室。\\
\hspace*{2em}上遂郊雍。至陇西。西登空桐。幸甘泉。令祠官宽舒等具泰一祠坛。坛放薄忌泰一坛。坛三垓。五帝坛环居其下。各如其方。黄帝西南。除八通鬼道。泰一所用。如雍一畤物。而加醴枣脯之属。杀一牦牛以为俎豆牢具。而五帝独有俎豆醴进。其下四方地。为餟食群神从者及北斗云。已祠。胙余皆燎之。其牛色白。鹿居其中。彘在鹿中。水而洎之。祭日以牛。祭月以羊彘特。泰一祝宰则衣紫及绣。五帝各如其色。日赤。月白。\\
\hspace*{2em}十一月辛巳朔旦冬至。昧爽。天子始郊拜泰一。朝朝日。夕夕月。则揖。而见泰一如雍礼。其赞飨曰。﹁天始以宝鼎神䇲授皇帝。朔而又朔。终而复始。皇帝敬拜见焉。﹂而衣上黄。其祠列火满坛。坛旁烹炊具。有司云。﹁祠上有光焉。﹂公卿言。﹁皇帝始郊见泰一云阳。有司奉瑄玉嘉牲荐飨。是夜有美光。及昼。黄气上属天。﹂太史公。祠官宽舒等曰。﹁神灵之休。祐福兆祥。宜因此地光域立泰畤坛以明应。令太祝领。秋及腊间祠。三岁天子一郊见。﹂\\
\hspace*{2em}其秋。为伐南越。告祷泰一。以牡荆画幡日月北斗登龙。以象天一三星。为泰一锋。名曰﹁灵旗﹂。为兵祷。则太史奉以指所伐国。而五利将军使不敢入海。之泰山祠。上使人微随验。实无所见。五利妄言见其师。其方尽。多不雠。上乃诛五利。\\
\hspace*{2em}其冬。公孙卿候神河南。见仙人迹缑氏城上。有物若雉。往来城上。天子亲幸缑氏城视迹。问卿。﹁得毋效文成。五利乎。﹂卿曰。﹁仙者非有求人主。人主求之。其道非少宽假。神不来。言神事。事如迂诞。积以岁乃可致。﹂于是郡国各除道。缮治宫观名山神祠所。以望幸矣。\\
\hspace*{2em}其年。既灭南越。上有嬖臣李延年以好音见。上善之。下公卿议。曰。﹁民间祠尚有鼓舞之乐。今郊祠而无乐。岂称乎。﹂公卿曰。﹁古者祀天地皆有乐。而神祇可得而礼。﹂或曰。﹁泰帝使素女鼓五十弦瑟。悲。帝禁不止。故破其瑟为二十五弦。﹂于是塞南越。祷祠泰一。后土。始用乐舞。益召歌儿。作二十五弦及箜篌瑟自此起。\\
\hspace*{2em}其来年冬。上议曰。﹁古者先振兵泽旅。然后封禅。﹂乃遂北巡朔方。勒兵十余万。还祭黄帝冢桥山。泽兵须如。上曰。﹁吾闻黄帝不死。今有冢。何也。﹂或对曰。﹁黄帝已仙上天。群臣葬其衣冠。﹂即至甘泉。为且用事泰山。先类祠泰一。\\
\hspace*{2em}自得宝鼎。上与公卿诸生议封禅。封禅用希旷绝。莫知其仪礼。而群儒采封禅尚书周官王制之望祀射牛事。齐人丁公年九十余。曰。﹁封者。合不死之名也。秦皇帝不得上封。陛下必欲上。稍上即无风雨。遂上封矣。﹂上于是乃令诸儒习射牛。草封禅仪。数年。至且行。天子既闻公孙卿及方士之言。黄帝以上封禅。皆致怪物与神通。欲放黄帝以尝接神仙人蓬莱士。高世比德于九皇。而颇采儒术以文之。群儒既以不能辩明封禅事。又牵拘于诗书古文而不敢骋。上为封祠器示群儒。群儒或曰﹁不与古同﹂。徐偃又曰﹁太常诸生行礼不如鲁善﹂。周霸属图封事。于是上绌偃。霸。尽罢诸儒弗用。\\
\hspace*{2em}三月。遂东幸缑氏。礼登中岳太室。从官在山下闻若有言﹁万岁﹂云。问上。上不言。问下。下不言。于是以三百户封太室奉祠。命曰崇高邑。东上泰山。山之草木叶未生。乃令人上石立之泰山颠。\\
\hspace*{2em}上遂东巡海上。行礼祠八神。齐人之上疏言神怪奇方者以万数。然无验者。乃益发船。令言海中神山者数千人求蓬莱神人。公孙卿持节常先行候名山。至东莱。言夜见一人。长数丈。就之则不见。见其迹甚大。类禽兽云。群臣有言见一老父牵狗。言﹁吾欲见巨公﹂。已忽不见。上既见大迹。未信。及群臣有言老父。则大以为仙人也。宿留海上。与方士传车及间使求仙人以千数。\\
\hspace*{2em}四月。还至奉高。上念诸儒及方士言封禅人人殊。不经。难施行。天子至梁父。礼祠地主。乙卯。令侍中儒者皮弁荐绅。射牛行事。封泰山下东方。如郊祠泰一之礼。封广丈二尺。高九尺。其下则有玉牒书。书秘。礼毕。天子独与侍中奉车子侯上泰山。亦有封。其事皆禁。明日。下阴道。丙辰。禅泰山下址东北肃然山。如祭后土礼。天子皆亲拜见。衣上黄而尽用乐焉。江淮间一茅三嵴为神藉。五色土益杂封。纵远方奇兽蜚禽及白雉诸物。颇以加祠。兕旄牛犀象之属弗用。皆至泰山然后去。封禅祠。其夜若有光。昼有白云起封中。\\
\hspace*{2em}天子从封禅还。坐明堂。群臣更上寿。于是制诏御史。﹁朕以眇眇之身承至尊。兢兢焉惧弗任。维德菲薄。不明于礼乐。脩祀泰一。若有象景光。屑如有望。依依震于怪物。欲止不敢。遂登封泰山。至于梁父。而后禅肃然。自新。嘉与士大夫更始。赐民百户牛一酒十石。加年八十孤寡布帛二匹。复博。奉高。蛇丘。历城。毋出今年租税。其赦天下。如乙卯赦令。行所过毋有复作。事在二年前。皆勿听治。﹂又下诏曰。﹁古者天子五载一巡狩。用事泰山。诸侯有朝宿地。其令诸侯各治邸泰山下。﹂\\
\hspace*{2em}天子既已封禅泰山。无风雨菑。而方士更言蓬莱诸神山若将可得。于是上欣然庶几遇之。乃复东至海上望。冀遇蓬莱焉。奉车子侯暴病。一日死。上乃遂去。并海上。北至碣石。巡自辽西。历北边至九原。五月。返至甘泉。有司言宝鼎出为元鼎。以今年为元封元年。\\
\hspace*{2em}其秋。有星茀于东井。后十余日。有星茀于三能。望气王朔言。﹁候独见其星出如瓠。食顷复入焉。﹂有司言曰。﹁陛下建汉家封禅。天其报德星云。﹂\\
\hspace*{2em}其来年冬。郊雍五帝。还。拜祝祠泰一。赞飨曰。﹁德星昭衍。厥维休祥。寿星仍出。渊耀光明。信星昭见。皇帝敬拜泰祝之飨。﹂\\
\hspace*{2em}其春。公孙卿言见神人东莱山。若云﹁见天子﹂。天子于是幸缑氏城。拜卿为中大夫。遂至东莱。宿留之数日。毋所见。见大人迹。复遣方士求神怪采芝药以千数。是岁旱。于是天子既出毋名。乃祷万里沙。过祠泰山。还至瓠子。自临塞决河。留二日。沉祠而去。使二卿将卒塞决河。河徙二渠。复禹之故迹焉。\\
\hspace*{2em}是时既灭南越。越人勇之乃言﹁越人俗信鬼。而其祠皆见鬼。数有效。昔东瓯王敬鬼。寿至百六十岁。后世谩怠。故衰秏﹂。乃令越巫立越祝祠。安台无坛。亦祠天神上帝百鬼。而以鸡卜。上信之。越祠鸡卜始用焉。\\
\hspace*{2em}公孙卿曰。﹁仙人可见。而上往常遽。以故不见。今陛下可为观。如缑氏城。置脯枣。神人宜可致。且仙人好楼居。﹂于是上令长安则作蜚廉桂观。甘泉则作益延寿观。使卿持节设具而候神人。乃作通天台。置祠具其下。将招来神仙之属。于是甘泉更置前殿。始广诸宫室。夏。有芝生殿防内中。天子为塞河。兴通天台。若有光云。乃下诏曰。﹁甘泉防生芝九茎。赦天下。毋有复作。﹂\\
\hspace*{2em}其明年。伐朝鲜。夏。旱。公孙卿曰。﹁黄帝时封则天旱。干封三年。﹂上乃下诏曰。﹁天旱。意干封乎。其令天下尊祠灵星焉。﹂\\
\hspace*{2em}其明年。上郊雍。通回中道。巡之。春。至鸣泽。从西河归。\\
\hspace*{2em}其明年冬。上巡南郡。至江陵而东。登礼潜之天柱山。号曰南岳。浮江。自寻阳出枞阳。过彭蠡。祀其名山川。北至琅邪。并海上。四月中。至奉高脩封焉。\\
\hspace*{2em}初。天子封泰山。泰山东北址古时有明堂处。处险不敞。上欲治明堂奉高旁。未晓其制度。济南人公玊带上黄帝时明堂图。明堂图中有一殿。四面无壁。以茅盖。通水。圜宫垣为复道。上有楼。从西南入。命曰昆仑。天子从之入。以拜祠上帝焉。于是上令奉高作明堂汶上。如带图。及五年脩封。则祠泰一。五帝于明堂上坐。令高皇帝祠坐对之。祠后土于下房。以二十太牢。天子从昆仑道入。始拜明堂如郊礼。礼毕。燎堂下。而上又上泰山。有秘祠其颠。而泰山下祠五帝。各如其方。黄帝并赤帝。而有司侍祠焉。泰山上举火。下悉应之。\\
\hspace*{2em}其后二岁。十一月甲子朔旦冬至。推历者以本统。天子亲至泰山。以十一月甲子朔旦冬至日祠上帝明堂。每脩封禅。其赞飨曰。﹁天增授皇帝泰元神策。周而复始。皇帝敬拜泰一。﹂东至海上。考入海及方士求神者。莫验。然益遣。冀遇之。\\
\hspace*{2em}十一月乙酉。柏梁灾。十二月甲午朔。上亲禅高里。祠后土。临渤海。将以望祠蓬莱之属。冀至殊庭焉。\\
\hspace*{2em}上还。以柏梁灾故。朝受计甘泉。公孙卿曰。﹁黄帝就青灵台。十二日烧。黄帝乃治明庭。明庭。甘泉也。﹂方士多言古帝王有都甘泉者。其后天子又朝诸侯甘泉。甘泉作诸侯邸。勇之乃曰。﹁越俗有火灾。复起屋必以大。用胜服之。﹂于是作建章宫。度为千门万户。前殿度高未央。其东则凤阙。高二十余丈。其西则唐中。数十里虎圈。其北治大池。渐台高二十余丈。名曰泰液池。中有蓬莱。方丈。瀛洲。壶梁。象海中神山龟鱼之属。其南有玉堂。璧门。大鸟之属。乃立神明台。井干楼。度五十余丈。辇道相属焉。\\
\hspace*{2em}夏。汉改历。以正月为岁首。而色上黄。官名更印章以五字。因为太初元年。是岁。西伐大宛。蝗大起。丁夫人。雒阳虞初等以方祠诅匈奴。大宛焉。\\
\hspace*{2em}其明年。有司言雍五畤无牢熟具。芬芳不备。乃命祠官进畤犊牢具。五色食所胜。而以木禺马代驹焉。独五帝用驹。行亲郊用驹。及诸名山川用驹者。悉以木禺马代。行过。乃用驹。他礼如故。\\
\hspace*{2em}其明年。东巡海上。考神仙之属。未有验者。方士有言﹁黄帝时为五城十二楼。以候神人于执期。命曰迎年﹂。上许作之如方。名曰明年。上亲礼祠上帝。衣上黄焉。\\
\hspace*{2em}公玊带曰。﹁黄帝时虽封泰山。然风后。封钜。岐伯令黄帝封东泰山。禅凡山合符。然后不死焉。﹂天子既令设祠具。至东泰山。东泰山卑小。不称其声。乃令祠官礼之。而不封禅焉。其后令带奉祠候神物。夏。遂还泰山。修五年之礼如前。而加禅祠石闾。石闾者。在泰山下址南方。方士多言此仙人之闾也。故上亲禅焉。\\
\hspace*{2em}其后五年。复至泰山脩封。还过祭常山。\\
\hspace*{2em}今天子所兴祠。泰一。后土。三年亲郊祠。建汉家封禅。五年一脩封。薄忌泰一及三一。冥羊。马行。赤星。五。宽舒之祠官以岁时致礼。凡六祠。皆太祝领之。至如八神诸神。明年。凡山他名祠。行过则祀。去则已。方士所兴祠。各自主。其人终则已。祠官弗主。他祠皆如其故。今上封禅。其后十二岁而还。遍于五岳。四渎矣。而方士之候祠神人。入海求蓬莱。终无有验。而公孙卿之候神者。犹以大人迹为解。无其效。天子益怠厌方士之怪迂语矣。然终羁縻弗绝。冀遇其真。自此之后。方士言祠神者弥众。然其效可睹矣。\\
\hspace*{2em}太史公曰。余从巡祭天地诸神名山川而封禅焉。入寿宫侍祠神语。究观方士祠官之言。于是退而论次自古以来用事于鬼神者。具见其表里。后有君子。得以览焉。至若俎豆圭币之详。献酬之礼。则有司存焉。

% ===== 索隐述赞（司马贞《史记索隐》）· 统一样式 =====
% 题头「索隐述赞曰」五字并入反白框；正文每四句一列，定宽盒 \hfil 居列中；
% \fontsize{24}{41} 列距对齐全书界行（121+41k）；\clearpage 使述赞独占一页。
\clearpage
\shuzanhead{索隐述赞曰}\linebreak
{\shuzanbody\fontsize{24}{41}\selectfont
\hbox to \hsize{\hfil 孝武纂极。四海承平。志尚奢丽。尤敬神明。\hfil}\linebreak
\hbox to \hsize{\hfil 坛开八道。接通五城。朝亲五利。夕拜文成。\hfil}\linebreak
\hbox to \hsize{\hfil 祭非祀典。巡乖卜征。登嵩勒岱。望景传声。\hfil}\linebreak
\hbox to \hsize{\hfil 迎年祀日。改历定正。疲耗中土。事彼边兵。\hfil}\linebreak
\hbox to \hsize{\hfil 日不暇给。人无聊生。俯观嬴政。几欲齐衡。\hfil}%
\par}
\end{document}
